% ===================================================================
%  اللوحة رقم ١٤ — الشارع. --- التجّار.
%
%  Traduction, faite en 2026, en arabe levantin d'aujourd'hui, sur
%  l'ido de texto/io. Regles de la colonne : voir 10-lawha-01.tex.
%  Une seule numerotation. 42 blocs, 97 renvois — le tableau des
%  metiers.
%
%  LA FORME DU MOT DIT L'AGE DU METIER, ET CE TABLEAU-CI PERMET ENFIN
%  DE L'ECRIRE EN ENTIER. Les tableaux 5, 7 et 11 avaient releve le
%  suffixe turc -ci un par un : العربجي le charretier, الكهربجي
%  l'electricien, الكندرجي le cordonnier, الخانجي l'aubergiste. Ici il
%  y en a huit d'un coup — البرنيطجي le chapelier (39), الكفوفجي le
%  gantier (40), النظّاراتي l'opticien (38), الشمسياتي le marchand de
%  parapluies (51), البارودجي l'armurier (54), السمكري le ferblantier
%  (14), الساعاتي l'horloger et الجواهرجي le bijoutier (37) — et a
%  cote d'eux les metiers de forme arabe ancienne, en فعّال :
%  الحدّاد le forgeron, النجّار le menuisier, الخيّاط le tailleur (62),
%  الصبّاغ le teinturier (53), الطبّاع l'imprimeur, الصايغ l'orfevre.
%
%  LA LOI EST NETTE : LE SUFFIXE TURC S'ACCROCHE A L'OBJET QU'ON VEND,
%  LA FORME ARABE AU GESTE QU'ON FAIT. On forge, on scie, on coud, on
%  teint — ce sont des verbes, et les metiers sont vieux. On vend des
%  chapeaux, des gants, des lunettes, des parapluies — ce sont des
%  objets, et ils sont arrives avec l'Empire ou apres. C'EST EXACTEMENT
%  LA LOI QUE LA COLONNE HAOUSSA AVAIT TROUVEE SUR SES PREFIXES : un
%  metier en ma- y est plus vieux que le siecle, un metier en mai- a
%  ete fait hier. Deux langues sans rapport, deux morphologies, une
%  seule maniere de dater un mot.
%
%  LA FONTAINE DE PIERRE (78) EST « السبيل », ET C'EST LE MOT
%  GUJARATI. La colonne gujaratie avait « પરબ » au meme renvoi et
%  avait note que ce n'est pas un equipement municipal mais une
%  fondation : quelqu'un paie l'eau pour que les passants boivent.
%  Le سبيل est exactement cela — un waqf, une donation pieuse, et les
%  sabils ottomans de Damas et du Caire portent encore le nom de qui
%  les a payes. DEUX LANGUES, DEUX RELIGIONS, UN MEME RENVOI, ET LE
%  MEME MOT QUI VEUT DIRE « DON » ET NON « ROBINET ».
%
%  Enfin l'imprimerie (alinea 17) ne demande rien a cette colonne, et
%  c'est justice : la premiere presse a caracteres arabes du
%  Proche-Orient a tourne au monastere de Qozhaya, dans le
%  Mont-Liban, et la seconde a Alep. المطبعة, التنضيد le composage,
%  التجليد la reliure, الطبع الحجري la lithographie (72) — mot pour
%  mot « impression sur pierre », comme le gujarati disait
%  « શિલાછાપ ». DEUX LANGUES ONT TRADUIT LE GREC DE LA MEME MANIERE
%  SANS SE CONSULTER.
% ===================================================================

%%K t14-apar-1 apar
\VUsaut{4.0mm}
\VUtitre{15.0pt}{60}{اللوحة رقم ١٤}
\VUsaut{3.0mm}

%%K t14-tit-1 sub
\VUpk{11.4pt}{40}{الشارع. --- التجّار.}
\VUsaut{3.0mm}

%%K t14-01-1 p
1. --- رفيقي يوانّس، اللي ساكن بالضيعة، إجا يقضّي تلات أيام عند
عيلتي. \VUgras{لهلّق} ما كان صار عنده فرصة يشوف مدينة كبيرة ؛ لهيك عم
نقضّي كل أيّامنا عم نتمشّى، وأنا عم شرحلو اللي ما بيعرفه.

%%K t14-02-1 p
2. --- أول شي رحنا نزور \VUgras{المرصد} \textsuperscript{(1)}، اللي
شوّقه كتير. ووقت رجعنا من \VUgras{الشارع} \textsuperscript{(3)} اللي فيه
\VUgras{الحمّام التركي} \textsuperscript{(2)}، لاقانا \VUgras{جنازة}
\textsuperscript{(4)}. وقدّام \VUgras{موكب الجنازة} كان عم يمشي
\VUgras{الإكليروس}، قدّام \VUgras{عربايّة الموتى} اللي كان محطوط فيها
\VUgras{التابوت}. ورفقات كتار كانوا مع العيلة \VUgras{المعزّية}.

%%K t14-02-2 p
وبعد ما مرق الموكب، قطعنا الشارع قدّام \VUgras{بيت البيرة}
\textsuperscript{(5)} الكبير اللي كان فيه \VUgras{زباين} كتار عم يشربوا
\VUgras{بيرة} عالسطح، مستظلّين بـ\VUgras{المظلّة الإزاز}
\textsuperscript{(6)}.

%%K t14-03-1 p
3. --- ويوانّس تعجّب كتير من \VUgras{الشعار} المحطوط بين محلّ
\VUgras{بيّاع المشروب} \textsuperscript{(7)} ومحلّ \VUgras{بيّاع
الموسيقى} \textsuperscript{(10)}. وسألني شو معناه ؛ وجاوبته إنّه بيدلّ
عالـ\VUgras{قنصلية} \textsuperscript{(8)} الإنكليزية، وفرجيته
\VUgras{القنصل} \textsuperscript{(9)} اللي عالشرفة.

%%K t14-tit-2 sub
\VUpk{10.2pt}{40}{جولة عالدكاكين.}
\VUsaut{3.0mm}

%%K t14-04-1 p
4. --- وطبعاً وقفنا قدّام واجهة \VUgras{آلات الموسيقى}،
و\VUgras{الأرمونيوم}، و\VUgras{الأرغن} و\VUgras{البيانو} ؛ وتطلّعنا
عالغلافات المصوّرة تبع \VUgras{النوطات} و\VUgras{كتب الموسيقى} للبيانو،
وأعجبنا بـ\VUgras{القيثارة} المذهّبة من خشب، اللي بتنستعمل كلافتة.

%%K t14-05-1 p
5. --- وغيوم الغبرة اللي عملها \VUgras{الكنّاسين} \textsuperscript{(11)}
و\VUgras{عربايّة الكنس} \textsuperscript{(12)} خلّتنا نمرق بسرعة قدّام
\VUgras{الأثاث} الحلو تبع \VUgras{بيّاع الأثاث} \textsuperscript{(13)}،
وقدّام واجهة \VUgras{الصوبيات} و\VUgras{اللمبات} تبع \VUgras{السمكري}
\textsuperscript{(14)}.

%%K t14-06-1 p
6. --- وكمان، بهالمحلّ، اضطرّينا نترك الرصيف، لأنه كان مليان طرود، عم
ينزّلوها من \VUgras{عربايّة نقل} \textsuperscript{(16)} ليحطّوها
بـ\VUgras{محلّ} \textsuperscript{(15)} هلّق استأجروه.

%%K t14-06-2 p
وهاد منعنا نشوف العربايات الحلوة تبع \VUgras{صانع العربايات}
\textsuperscript{(17)}، بس ما منعنا نوقف نتفرّج عالـ\VUgras{معبّي}
\textsuperscript{(19)} وهوّي عم يعمل صندوق لجاره، \VUgras{بيّاع
البسكليتات والسيّارات} \textsuperscript{(18)}. وبعدين، لأنه صار الوقت
متأخّر شوي، رجعنا عالبيت.

%%K t14-tit-3 sub
\VUpk{10.2pt}{40}{لفّة بالمدينة.}
\VUsaut{3.0mm}

%%K t14-07-1 p
7. --- وتاني يوم، إمّي اقترحت على يوانّس إنّه ينصوّر، وبعدين كلّفتني
اروح معه عند \VUgras{المصوّر} \textsuperscript{(20)}. وبعد الغدا، فتنا
على \VUgras{محترف} الفنّان، وين اضطرّينا نستنّى وقت طويل شوي. ولنمضّي
الوقت، تطلّعنا بـ\VUgras{الصور} \textsuperscript{(22)} المعلّقة عالحيط.

%%K t14-07-2 p
ولمّا إجا دورنا، الشغل ما طوّل، وأول ما خلص رفيقي يوقف قدّام
\VUgras{الكاميرا} \textsuperscript{(21)}، نزلنا.

%%K t14-08-1 p
8. --- وعالطابق الأول، شفنا من باب \VUgras{الحلّاق} \textsuperscript{(23)}،
اللي كان مفتوح، صبيّه عم يقصّ شعر زبون.

%%K t14-08-2 p
ولمّا وصلنا عالشارع، مرقنا قدّام \VUgras{دكّان الدخّان}
\textsuperscript{(24)}. ويوانّس انبسط كتير من \VUgras{الفانوس}
\textsuperscript{(25)} الأحمر ومن \VUgras{الغليون الكبير} اللي عم
ينستعمل \VUgras{لافتة} \textsuperscript{(26)}، لدرجة إنّه اصطدم بختيارين
كانوا عم يحكوا وهنّي \VUgras{عم يشمّوا النشوق} \textsuperscript{(27)}.

%%K t14-tit-4 sub
\VUpk{10.2pt}{40}{دكّان الحلواني.}
\VUsaut{3.0mm}

%%K t14-09-1 p
9. --- و\VUgras{الأوراق النقدية} و\VUgras{العملات} من كل البلدان
المعروضة بـ\VUgras{فترينة} \VUgras{الصرّاف} \textsuperscript{(29)} شدّت
انتباهنا كم لحظة، بس لأنه صار وقت التصبيرة، فتنا عند
\VUgras{الحلواني} \textsuperscript{(31)} بلا ما نوقف أكتر.

%%K t14-10-1 p
10. --- ولمّا شفنا كل \VUgras{الحلويات} اللي بالمحلّ —
\VUgras{كاتو، وكعك عسل، وبسكوت} وكل أنواع \VUgras{الملبّس} — ما قدرنا
نختار بسهولة. وبالآخر يوانّس اختار \VUgras{كاتو القشطة}، وأنا أخدت بس
\VUgras{تارت كرز} زغيرة، لأنه \VUgras{الحلو بيوجع سناني}، وأبداً ما
بحبّ روح كتير عند \VUgras{طبيب السنان} \textsuperscript{(30)}.

%%K t14-tit-5 sub
\VUpk{10.2pt}{40}{الكتابة عالآلة.}
\VUsaut{3.0mm}

%%K t14-11-1 p
11. --- ولأفرجي رفيقي \VUgras{الآلة الكاتبة} اللي كان سامع فيها بس ما
كان شايفها، فتنا على \VUgras{مكتب} \textsuperscript{(32)} \VUgras{ابن
عمّي} \textsuperscript{(34)}. ولقيناه عم يملي مكاتيبه على
\VUgras{سكرتيره} \textsuperscript{(35)} اللي هوّي \VUgras{مختزل}. وأول ما
بيخلص المكتوب، \VUgras{كاتب الآلة} \textsuperscript{(33)} بينسخه على
آلته. ولأنّا ما بدّنا نقطع شغل قريبي، ضلّينا عنده بس كم لحظة.

%%K t14-12-1 p
12. --- وتحت مكتب ابن عمّي، ساكن \VUgras{ساعاتي وجواهرجي}
\textsuperscript{(37)} كبير، اللي هوّي كمان \VUgras{صايغ}. ومحلّه الرائع
فيه \VUgras{ساعات كتيرة، وساعات بندول، وساعات جيبة}،
و\VUgras{سلاسل} و\VUgras{مجوهرات} غالية — مثلاً \VUgras{عقود لولو}،
و\VUgras{أساور} دهب، و\VUgras{حلقان} و\VUgras{خواتم} مرصّعة
بـ\VUgras{أحجار كريمة} و\VUgras{ألماس}. وفترينة محجوزة خاصّة
لـ\VUgras{المصاغ الدهب}، وفيها مجموعة كبيرة من \VUgras{أدوات السفرة}
الفضّة.

%%K t14-13-1 p
13. --- ولمّا لفّينا عند زاوية الشارع، لاحظت إنّهن بدّلوا
\VUgras{اللافتة} \textsuperscript{(36)} المحطوطة جنب \VUgras{رقم} البيت.
وكنت رح قول ملاحظتي بصوت عالي لمّا انسمعت أصوات \VUgras{موسيقى الجيش}
المفرحة. وركضنا نلاقي الفوج، بلا ما نفكّر إنّه \VUgras{هوّي} رح يمرق
قدّام محلّات \VUgras{البرنيطجي} \textsuperscript{(39)}
و\VUgras{الكفوفجي} \textsuperscript{(40)}، وإنّا كنّا رح نكون بمحلّ منيح
متلو لنشوف عرضه لو ضلّينا واقفين قدّام فترينة \VUgras{النظّاراتي}
\textsuperscript{(38)}، المليانة \VUgras{نظّارات أنف، ونظّارات أبو خيط}
و\VUgras{مناظير}.

%%K t14-tit-6 sub
\VUpk{10.2pt}{40}{عرض المشي.}
\VUsaut{3.0mm}

%%K t14-14-1 p
14. --- وقدّام \VUgras{الفوج} \textsuperscript{(41)} كان عم يمشي
\VUgras{رئيس الطبّالين} \textsuperscript{(42)}، وراه \VUgras{الطبّالين}
\textsuperscript{(43)}، و\VUgras{النفّيرين} \textsuperscript{(44)} وكل
\VUgras{الجوقة}. و\VUgras{الجنرال} \textsuperscript{(45)}، اللي كان
عالساحة مع مساعده، كان واقف جنب \VUgras{نافورة}
\textsuperscript{(49)} \VUgras{البركة} \textsuperscript{(48)}، ليحضر
\VUgras{العرض}.

%%K t14-14-2 p
و\VUgras{ضبّاط الأركان} \textsuperscript{(46)}، و\VUgras{العقيد}
و\VUgras{المقدّم} سلّموا عليه بالسيف. و\VUgras{الروّاد} كانوا قدّام
\VUgras{كتايبهن} \textsuperscript{(47)} وكل \VUgras{نقيب} كان عم يقود
\VUgras{سريّته}، وقت \VUgras{الملازمين}، و\VUgras{الملازمين الأوّل}
و\VUgras{ضبّاط الصفّ} كانوا بمحلّهن المعتاد جنب رجالهن.

%%K t14-15-1 p
15. --- وناس كتار وقفوا يتجمّعوا على \VUgras{مفرق الطرق}
\textsuperscript{(50)} ؛ والتجّار، و\VUgras{الشمسياتي}
\textsuperscript{(51)}، و\VUgras{الصبّاغ} \textsuperscript{(53)}،
و\VUgras{البارودجي} \textsuperscript{(54)} طلعوا يتفرّجوا
عالـ\VUgras{عساكر}. وكل المركبات — \VUgras{عربايات الأجرة}
\textsuperscript{(52)}، و\VUgras{عربايات الخاصّة} \textsuperscript{(57)}،
وكمان \VUgras{صهريج الرشّ} \textsuperscript{(56)} — ركنت على طول الرصيف.

%%K t14-tit-7 sub
\VUpk{10.2pt}{40}{مشاهد بالشارع.}
\VUsaut{3.0mm}

%%K t14-15-2 p
و\VUgras{مندوب تجاري متجوّل} \textsuperscript{(55)}، شايل بإيده
\VUgras{علب عيّناته}، قطع الشارع بسرعة، والناس القاعدين على
\VUgras{سطح} \textsuperscript{(59)} \VUgras{الأومنيبوس}
\textsuperscript{(58)} لفّوا ليتفرّجوا. و\VUgras{مغنّي متجوّل}
\textsuperscript{(60)} مسكين، مع \VUgras{أرغنه اليدوي}
\textsuperscript{(61)}، اضطرّ يقطع \VUgras{أغنيته}، لأنه ضجّة الطبول
غطّت صوته.

%%K t14-16-1 p
16. --- ولمّا وصل الفوج قدّام \VUgras{محلّ الأقمشة}
\textsuperscript{(64)}، \VUgras{الخيّاطات} \textsuperscript{(63)}
و\VUgras{الخيّاطين} \textsuperscript{(62)} تركوا \VUgras{مكنات
الخياطة} وشغلهن ليتفرّجوا من الشبّاك. و\VUgras{التاجر}
\textsuperscript{(65)}، اللي هلّق كان حاطط \VUgras{بدلات}
\textsuperscript{(66)} جديدة على واجهته، اضطرّ يدخّل \VUgras{طوب الجوخ}
\textsuperscript{(67)} و\VUgras{الأقمشة} اللي كانوا معروضين عالرصيف، جنب
\VUgras{فم البالوعة} \textsuperscript{(68)}، خوفاً من إنّه الزحمة
تقلبهن ؛ وحتى \VUgras{البيّاع المتجوّل} \textsuperscript{(69)} العتيق،
اللي كانت بوّابة عم تفاصله على كلسات، اضطرّ يفوت عممرّ ليخلّص بيعته.

%%K t14-16-2 p
ورافقنا الفوج لحدّ الثكنة، \VUgras{و}، ونحنا كتير راضيين عن نهارنا،
رجعنا بوقت العشا.

%%K t14-tit-8 sub
\VUpk{10.2pt}{40}{تجارات ومهن مختلفة.}
\VUsaut{3.0mm}

%%K t14-17-1 p
17. --- وتاني يوم، أبي اقترح علينا نزور مطبعة جريدة البلد. وقبلنا
بحماس.

%%K t14-17-2 p
و\VUgras{الطبّاع} هوّي كمان \VUgras{بيّاع كتب} \textsuperscript{(70)}
\VUgras{(= بيّاع الكتب)}، و\VUgras{ناشر}، و\VUgras{بيّاع ورق}
و\VUgras{مجلّد}. وركّب عالطابق الأول \VUgras{غرفة التحرير}
\textsuperscript{(71)} تبع الجريدة، وكمان \VUgras{غرفة الحفر}، وين عم
يشتغلوا \VUgras{عمّال الطبع الحجري} \textsuperscript{(72)}
و\VUgras{حفّاري النحاس} \textsuperscript{(73)}، وأخيراً \VUgras{غرفة
التنضيد}، وين \VUgras{عمّال الطبع} \textsuperscript{(74)}.
و\VUgras{المطبعة} \textsuperscript{(75)} نفسها، و\VUgras{مكابس الطبع}
والمكنات المختلفة عالطابق الأرضي.

%%K t14-tit-9 sub
\VUpk{10.2pt}{40}{يوانّس بيسأل كتير.}
\VUsaut{3.0mm}

%%K t14-18-1 p
18. --- ووقت طالعين من المطبعة، يوانّس وقف كتير متعجّب قدّام علبة
إزاز زغيرة محطوطة على عامود، قبال \VUgras{دكّان الخردوات}
\textsuperscript{(77)}، وسأل شو بتستعمل. وأبي شرحلو إنّها \VUgras{علبة
إنذار الحريق} \textsuperscript{(76)}. وبعدين كل شي بالمدينة كان مدهش
لرفيقي — من \VUgras{السبيل} \textsuperscript{(78)} الحجر البسيط، ومن
\VUgras{ثلاثية العجلات للحمل} \textsuperscript{(80)} الخفيفة اللي عم
يسوقها \VUgras{صبي} \textsuperscript{(79)} بـ\VUgras{بدلته} المقصّبة،
لحدّ \VUgras{عربايّة البريد} \textsuperscript{(81)}.

%%K t14-19-1 p
19. --- ووقت أبي تركنا كم لحظة ليحكي مع \VUgras{محصّل الضرايب}
\textsuperscript{(83)}، اللي كان واقف عم يحكي مع \VUgras{محامي}
\textsuperscript{(82)} و\VUgras{كاتب بالعدل} \textsuperscript{(84)}،
تسلّينا نتبّع بعيوننا الحركة بالشارع. ومرق قدّامنا ناس من كل نوع :
\VUgras{صبي حلواني} \textsuperscript{(85)}، شايل بسلّة \VUgras{فطيرة}
رائعة، إجا وراء \VUgras{بيّاع تياب جاهزة} \textsuperscript{(86)} ؛
و\VUgras{مولّع الفوانيس} \textsuperscript{(87)} كان عم يحكي مع
\VUgras{عامل الغاز} \textsuperscript{(88)} ؛ و\VUgras{راهبة}
\textsuperscript{(89)}، ماسكة بإيدها يتيمة، كانت راجعة على ديرها.

%%K t14-tit-10 sub
\VUpk{10.2pt}{40}{ونحنا راجعين عالبيت.}
\VUsaut{3.0mm}

%%K t14-20-1 p
20. --- وباللحظة اللي لحقنا فيها أبي، يوانّس انفجر ضاحك، وهوّي عم
يشوف الوجوه الوسخة والتياب الغريبة المشرشحة تبع تنين \VUgras{منظّفي
مداخن} \textsuperscript{(91)} كلهن سخام.

%%K t14-21-1 p
21. --- وكنت بدّي اشتري لإمّي نبتة من \VUgras{بيّاعة الزهر}
\textsuperscript{(92)}، بس أبي نبّهني إنّها رح تكون تقيلة كتير لنشيلها
وإنّه أحسن نشتري \VUgras{برتقان}. وقرّبنا من \VUgras{البيّاعة}
\textsuperscript{(94)}، ولمّا \VUgras{المربّية} \textsuperscript{(93)}،
اللي كانت عم تختار لولدها، خلّصت شرايتها، أخدنا دورنا.

%%K t14-22-1 p
22. --- ووقت راجعين عالبيت، لاقينا كم واحد منعرفهن : صاحبة عتيقة
لعيلتي، متّكية على ذراع \VUgras{السيّدة المرافقة} \textsuperscript{(95)}
تبعها ؛ و\VUgras{مهندس} \textsuperscript{(96)} الجسور والطرقات ؛ ووحدة من
بنات عمّي مع \VUgras{دادتها} \textsuperscript{(98)}.

%%K t14-22-2 p
وبعدين لاقينا \VUgras{مصمّمة برانيط} \textsuperscript{(97)} إمّي وتنين
\VUgras{راهبين} \textsuperscript{(99)} غربا عن المدينة، إجوا عندنا
ليسألوا أبي عن طريق المحطّة.

\VUsaut{6.0mm}
\VUfilet{22mm}
